\documentclass[english, a4paper, numbers=noenddot]{article}

\usepackage[utf8]{inputenc}
\usepackage[T1]{fontenc}
\usepackage[english]{babel}
\usepackage{lmodern}
\usepackage[top=1.7cm, right=1.8cm, bottom=1.8cm, left=1.8cm]{geometry}
\usepackage{calc}
\usepackage{graphicx}
\usepackage{amsmath}
\usepackage{tabularx}
\usepackage{booktabs}
\usepackage[official]{eurosym}
\usepackage{lastpage}
\usepackage{subcaption}
\usepackage{etoolbox}
\usepackage{amssymb}
\usepackage{enumitem}
\usepackage{natbib}
\usepackage{siunitx}
\usepackage[hang,flushmargin]{footmisc}
\usepackage{float} % for [H] figure placement
\usepackage[export]{adjustbox}

% spacing
\usepackage{caption}
\captionsetup{font=small}

\captionsetup[figure]{skip=5pt}
\captionsetup[subfigure]{font=scriptsize, skip=1pt}
\captionsetup[table]{skip=6pt}

\setlength{\textfloatsep}{9pt plus 1pt minus 2pt}
\setlength{\floatsep}{6pt plus 1pt minus 2pt}
\setlength{\intextsep}{8pt plus 1pt minus 2pt}
\setlength{\abovecaptionskip}{3pt}
\setlength{\belowcaptionskip}{0pt}

\bibliographystyle{plainnat}

% TikZ setup
\usepackage{tikz}
\usetikzlibrary{arrows.meta, positioning, shapes, calc, fit}

% Line spacing
\usepackage{setspace}
\setstretch{1.15}

% Page header
\usepackage[automark]{scrlayer-scrpage}
\pagestyle{scrheadings}
\clearscrheadings
\ihead{}
\chead{}
\ohead{\pagemark}

% Hyperref (should come last)
\usepackage[
  colorlinks=false,
  pdfborder={0 0 0}
]{hyperref}


% \setlength{\parindent}{0pt}
\usepackage{parskip}
\usepackage{wrapfig}

%-------------- Ende allgemeine Angaben -----------------------------
\begin{document}





\section{Socio-Economic Analysis}

\subsection{Aggregation of predicted PM2.5 to Kreis level}

The dense-grid predictions described in Section~\ref{sec:results} are produced
at regularly spaced locations rather than directly at administrative units.
To combine these predictions with the socioeconomic data, each grid point is
therefore assigned to the Kreis polygon containing its centre coordinates.
Predicted PM2.5 is then averaged across all grid points located within each
Kreis.

Of 1,728 usable grid predictions, 1,719 could be assigned to a Kreis. This
results in Kreis-level modeled pollution estimates for 134 Kreise within the
North-Eastern German prediction domain. The number of contributing grid
points varies with Kreis size, ranging from 1 to 54, with a median of 12 grid
points per Kreis.

Our main exposure measure is the arithmetic mean of the predicted PM2.5
values within each Kreis. It should therefore be interpreted as an
area-oriented estimate of mean modeled ambient pollution rather than as a
population-weighted measure of individual exposure. Constructing the latter
would require information on the within-Kreis spatial distribution of the
population.

Across the 134 covered Kreise, mean predicted PM2.5 ranges from approximately
6.02 to 13.12~\si{\micro\gram\per\cubic\meter}, with an overall Kreis mean of
9.88~\si{\micro\gram\per\cubic\meter}.

\subsection{Spatial patterns}

Figure~\ref{fig:socioeconomic_maps} compares the spatial distribution of mean
predicted PM2.5 with selected socioeconomic and demographic indicators.
All panels refer to the same modeled study region. Kreise outside the dense
prediction grid are shown only as geographic context and are not included in
the analysis.

\begin{figure}[ht]
    \centering
    \includegraphics[width=\textwidth]
    {Figures/socioeconomic/00_summary_maps.png}
    \caption{Mean predicted PM2.5 and selected socioeconomic and demographic
    characteristics across Kreise covered by the dense prediction grid.
    Population density is displayed using a logarithmic color scale.}
    \label{fig:socioeconomic_maps}
\end{figure}

The maps show substantial spatial variation in both modeled PM2.5 and the
socioeconomic indicators. However, the spatial patterns do not indicate a
single, consistent socioeconomic gradient in modeled pollution. Areas of
higher predicted PM2.5 overlap with some relatively dense and urbanized
Kreise, but similar pollution levels also occur across Kreise with different
income, unemployment, education, and demographic profiles. We therefore
evaluate these relationships quantitatively rather than inferring an exposure
gradient from the maps alone.

\subsection{Association with socioeconomic characteristics}

Table~\ref{tab:pm25_socioeconomic_correlations} reports descriptive Pearson
correlations between mean predicted PM2.5 and the selected Kreis-level
indicators. The associations are generally weak. The strongest positive
association is observed for unemployment ($r=0.225$), followed by university
degree share ($r=0.198$) and log population density ($r=0.171$). Disposable
income is weakly negatively associated with predicted PM2.5 ($r=-0.145$).
Immigration history shows essentially no linear relationship with predicted
PM2.5 ($r=-0.022$).

\begin{table}[ht]
\centering
\caption{Pearson correlations between Kreis-level mean predicted PM2.5 and
selected socioeconomic and demographic indicators. Correlations are
descriptive and do not imply causal relationships.}
\label{tab:pm25_socioeconomic_correlations}
\begin{tabular}{lcc}
\toprule
\textbf{Indicator} & \textbf{$n$} & \textbf{Pearson $r$} \\
\midrule
Disposable income per capita      & 132 & -0.145 \\
Unemployment rate                 & 134 &  0.225 \\
No vocational qualification       & 134 & -0.153 \\
University degree                 & 134 &  0.198 \\
Immigration history               & 134 & -0.022 \\
Log population density            & 134 &  0.171 \\
Population aged 65+               & 134 & -0.120 \\
\bottomrule
\end{tabular}
\end{table}

Overall, these results provide little evidence for a strong or uniform
socioeconomic gradient in modeled PM2.5 across the study region. The weak
negative association with disposable income and weak positive association
with unemployment are consistent with somewhat higher modeled pollution in
more economically disadvantaged Kreise, but their magnitude is small.
Moreover, other indicators do not follow the same pattern: the share without
a vocational qualification is weakly negatively associated with predicted
PM2.5, while university-degree share is weakly positively associated with it.
The near-zero association with immigration history further illustrates that
the observed spatial patterns cannot be summarized by a single dimension of
social disadvantage.

These results should also be interpreted in light of the limitations of the
pollution estimates themselves. The socioeconomic analysis is based on model
predictions rather than direct measurements, and Section~\ref{sec:results}
shows systematic overprediction and heterogeneous predictive performance in
the held-out German region. In addition, the Kreis-level PM2.5 measure is an
unweighted spatial average and does not account for where people live within
each Kreis. Consequently, the analysis should be understood as an exploratory
assessment of spatial associations between modeled ambient pollution and
Kreis-level socioeconomic characteristics, rather than as evidence of causal
environmental inequality or individual exposure differences.

https://www.eea.europa.eu/en/analysis/maps-and-charts/air-quality-statistics-dashboards

\section{Conclusion}

In this project, we build a model that predicts PM 2.5 pollution given exclusively publicly available satellite data. For a given location, patches are collected and the output is a PM2.5 point estimate. The results show better results than the baseline, suggesting that there is a slight tendency to overestimate pollution in the test region (NE Germany). Moreover, we provide evidence of the unreliability of low cost sensors and resort to using EEA-wide reference stations as labels for our models. We show that from scratch models perform better in estimating PM2.5 in the test region compared to the ResNet-50 model pre-trained on BigEarthNet. We offer possible explanations to why that is the case. Finally, we show that PM 2.5 correlates to several social indicators, underscoring that ......

There are several limitations to our project. 

The performance of the model could be further improved by including additional scalar variables, especially the ones related to metereological phenomena.

%how good is the model at predicting pollution 

%conclusion from socio economic analysis 

%implications: include the question is it ok to act on predicted no measured values

%limitations of our approach 

%Limitations 
%improvement of models, inclusion of additional covariates. 
%- regulations (factors with filters, regulation for cars) 
%- gdp 
%- power type used in region 
%- additional scalar values 
%- airflow 
%- larger spatial context (po-area) 

%future work?

\bibliographystyle{plain}
\bibliography{references}
\cite{who2021airquality}
\cite{scheibenreif2021airpollution,benabbas2019particulate}
\clearpage
\section*{Appendix}
\appendix

\section{Additional Graphs and Images}

\begin{figure}[ht]
    \centering
    
\includegraphics[width=0.2\textwidth]{Figures/sds011_sensor.jpg}
    \caption{SDS011 sensor, the most common low-cost sensor in Germany}
    \label{fig:sds011_sensor}
\end{figure}












\end{document}